
% ============================================================================
\clearpage
\section*{v4.9.0-r80 Source-Only: The Background of the OP5 Comparison System Is Not Sourced by Its Spinors; Closing the Comparison-Model Branch}
\addcontentsline{toc}{section}{v4.9.0-r80 Source-Only: The Background of the OP5 Comparison System Is Not Sourced by Its Spinors; Closing the Comparison-Model Branch}
\label{sec:v490-r80-background-not-sourced}
% ============================================================================

\paragraph{Purpose.}
Layer v4.9.0-r79 left one source of Lorentz breaking in the comparison system: the prescribed background
\((\Psi\text{ at the germ},A,h)\) (\S\ref{sec:v490-r79-gauge-independence}).  This layer asks whether the background is
generated by the three spinor blocks themselves.  It records one exploration (polar/Fierz variables of the blocks) and one
pre-registered test (H68), which is refuted.  On 4 November 2026 the author decided to close the comparison-model branch at
this point and to return to the common question behind OP1--OP5, a joint dynamics.  The layer therefore ends with a closing
register of layers r71--r80.  Everything concerns the comparison model of \S\ref{sec:v410-state-mediator} only.  Earlier
layers are not rewritten.

\noindent\textbf{Source-only status.}
This v4.9.0-r80 layer is source-only.  It does not build a PDF, does not assign a new DOI,
and does not authorize a public release.  The v4.9.0 Version DOI
\href{https://doi.org/10.5281/zenodo.23534954}{10.5281/zenodo.23534954} and the Concept DOI
\href{https://doi.org/10.5281/zenodo.19210728}{10.5281/zenodo.19210728} are unchanged by it.

\subsection*{v4.9.0-r80 results}
\label{subsec:v490-r80-results}

\paragraph{Theoremlet v4.9.0-r80.1 (exploration: the three blocks are regular spinor fluids in a Hubble-like flow).}
Exploration without a criterion.  At the germ \((t,x,y,z)=(1,0.6,0,0)\) the polar (Fierz) variables of each block
\(u_n\in\mathbb C^4\) were computed in the Dirac basis, with \(\Phi=\bar u u\), \(\Theta=\bar u i\gamma_5u\),
\(J^a=\bar u\gamma^a u\), \(K^a=\bar u\gamma^a\gamma_5u\).
\begin{itemize}[leftmargin=2em]
\item The Fierz identities \(J^2=\Phi^2+\Theta^2=-K^2\), \(J\cdot K=0\) hold to \(10^{-14}\).
\item All three blocks lie in the regular Dirac sector with Yvon--Takabayasi angle \(|\beta|\le0.07^\circ\); the densities
are \(\phi^2=5.97,\ 0.43,\ 6.38\).
\item The velocities are \(u^a=\gamma(1,\pm0.6,0,0)\), \(\gamma=1.25\): blocks 0 and 1 move with \(v=+x/t\), block 2 against
it (the counterflow of the earlier series).
\item All spin axes point along the azimuthal direction \(\hat\varphi=e_y\), parallel for blocks 0, 1 and antiparallel for
block 2.  The same axis is singled out by the gauge \(K_2\) (stabilizer \(\mathfrak{so}(1,2)\), r78.1) and by the background
\(A=M(t)\,d\varphi\).
\end{itemize}

\paragraph{Theoremlet v4.9.0-r80.2 (H68 refuted: no Yang--Mills or Klein--Gordon equation with spinor sources).}
Pre-registered.  The background is \(A=M(t)\,d\varphi\) with
\(M(t)=\tfrac{3}{2\cos2\theta}(\cos\theta\,\lambda_4-\sin\theta\,\lambda_6)\), \(\theta=\tfrac12\arcsin(t/2)\), and
\(h=\gamma/\cos2\theta\).  Known in advance: \(F=M'(t)\,dt\wedge d\varphi\), and the values of \(A,F\) generate
\(\mathfrak k=\mathrm{span}(\lambda_2,\lambda_4,\lambda_6)\cong\mathfrak{su}(2)\), with \(Y_3=\lambda_2\in\mathfrak k\).  In full
\(\mathfrak{su}(3)\) the test fails trivially, because the blocks carry a \(\lambda_3\) density; it was therefore restricted to
\(\mathfrak k\).
\begin{itemize}[leftmargin=2em]
\item \emph{Setup.}  Metric of the model's own coframe \(g=E^{\mathsf T}\eta E\) with \(E^a_\mu=\mathrm{Re}(\Psi^\dagger Q_aD_\mu\Psi)\),
\(\sqrt{|g|}=|\det E|\); \(D=\partial+iA\),
\(L^\nu=|g|^{-1/2}\partial_\mu(|g|^{1/2}F^{\mu\nu})+i[A_\mu,F^{\mu\nu}]\); colour current
\(j^\nu_b=\bar\Psi\gamma^\nu(\lambda_b\otimes1)\Psi\); scalar source \(s=\bar\Psi(Y\otimes1)\Psi\), the derivative of the Yukawa
term.  Truncated power series around the germ (repaired 4-jet), equations to degree 2, constants by least squares.
\item \emph{Hypotheses.}  (a) \(L^\nu_b=\kappa_Aj^\nu_b\) for \(b\in\{2,4,6\}\); (b) \(\Box_gh=\kappa_hs\); (c) \(\Box_gh+m^2h=\kappa_hs\).
Confirmed at relative residual \(r\le10^{-6}\), refuted at \(r\ge10^{-2}\).
\item \emph{Controls.}  Bianchi identity of \(L\) to degree 1: \(1.2\times10^{-14}\); flat \(\Box h\) from the series equals the
exact value \(-6.344838895781825\) to \(2\times10^{-15}\); germ residual of \(R\) \(\le2.8\times10^{-13}\).
\item \emph{Result.}  (a) \(r=0.955\) (per degree \(0.965,\ 0.963,\ 0.954\), fitted \(\kappa_A\) between \(0.12\) and \(22\));
(b) \(r=0.995\); (c) \(r=0.826\).  All three refuted.  Already degree 0, which does not depend on the higher jets and their
\(\vartheta\)-gauge, fails.
\item \emph{Why.}  At the germ the sources in \(\mathfrak k\) are about \(10^{-3}\), the left side is \(O(0.1\text{--}0.4)\); for
the scalar \(s=0.002\) against \(\Box_gh=6.91\).  The directions do not match either.  The spinor currents lie almost entirely
outside \(\mathfrak k\), in \(\lambda_0,\lambda_3,\lambda_8\) and \(\lambda_1,\lambda_5,\lambda_7\): the blocks are practically
incoherent with respect to the off-diagonal generators that carry \(A\).
\item \emph{Flat metric (exploration).}  The analytic structure is confirmed: \(L\) has only the components \((t,\lambda_2)\) and
\((\hat\varphi,\lambda_4,\lambda_6)\).  The fits fail as well (\(r=0.9997\), \(1.000\), \(0.28\)).
\end{itemize}

\paragraph{Theoremlet v4.9.0-r80.3 (closing register of the comparison-model branch, r71--r80).}
Author decision of 4 November 2026.  The comparison system is a fixed construction: its background is prescribed and, by
r80.2, not generated by its spinors through Yang--Mills or Klein--Gordon equations.  Its structural results are local and
concern this construction only.
\begin{itemize}[leftmargin=2em]
\item \emph{Established at the germ.}  An exact \(\vartheta\)-gauge symmetry (r71); involution of the gauged symbol from order
6 (r72); exactly four hidden conditions of order 6, supported on the light cone (r73, r77); formal integrability of the
extended system (H59); six real polarizations per null direction (H60); covariance of \(R\) and \(\Omega\) (H64, H66); gauge
independence of the four conditions (H67).
\item \emph{Not established.}  A native origin of the background (r80.2); any identification of the six modes with known
fields (H61); any statement about the QSC model itself; the OP5 vacuum term.
\item \emph{Closed as a line of work.}  Further refinements of this comparison system (Lorentz action on the six modes, second
base point) are not pursued without a native dynamics.  The next step is one level up: a survey of the action ansätze tried
so far and of what a joint action for \(\Psi\), connection, scalar and coframe must at least achieve.
\end{itemize}

\paragraph{Proof and measurement sources.}
r80.1 by \texttt{k168\_fierz.py} and the note \texttt{Erkundung\_\allowbreak Fierz\_\allowbreak Hintergrund.md}; r80.2 by
\texttt{k169\_\allowbreak feldgleichungen.py} with output \texttt{k169\_out.json}.  The note \texttt{Vorab\_\allowbreak Notiz\_\allowbreak H68}
and the result note lie in \texttt{notes/\allowbreak research/\allowbreak predictions/}.  \(\square\)

\paragraph{Consequence.}
The comparison system leaves no internal source of its Lorentz breaking: the breaking comes from a background that the
system does not generate.  The branch is closed with this layer.  OP5 is not closed; the open problem is moved to the level
of a joint dynamics, which OP1--OP5 share.

\subsection*{v4.9.0-r80 status in the three frames}
\label{subsec:v490-r80-frames}

\begin{itemize}[leftmargin=2em]
\item All statements concern the fixed OP5 comparison model in one chart, with one coframe reconstructed from \(\Psi\).  No
frame map M, T or P acts on them.
\item The colour labels \(\lambda_b\) of r80.2 are the internal \(3\times3\) indices of the comparison model, read in its
single fixed frame.  They are not identified with \(\mathrm{SU}(3)_c\) of \S\ref{sec:v490-r35-inner-outer-unit}.
\end{itemize}

\subsection*{v4.9.0-r80 register entries}
\label{subsec:v490-r80-register}

\begin{center}
\small
\begin{tabular}{@{}>{\raggedright\arraybackslash}p{0.34\textwidth}>{\raggedright\arraybackslash}p{0.58\textwidth}@{}}
\toprule
statement & status after v4.9.0-r80\\
\midrule
spinor blocks at the germ & regular Dirac sector, Hubble-like counterflow, spins along \(\hat\varphi\) (exploration)\\
background \(A,h\) & not sourced by the blocks via Yang--Mills on \(\mathfrak k\) or Klein--Gordon (H68 refuted)\\
comparison-model branch & closed (author decision); results r71--r79 remain valid as local statements\\
next step & survey of action ansätze; joint dynamics\\
\bottomrule
\end{tabular}
\end{center}

\subsection*{v4.9.0-r80 ledger and guard update}
\label{subsec:v490-r80-ledger-guard}

\[
R_{\rm chk}^{v490\text{-}r80}=(1,0,0)
\]
for (H68 decided; background not native; OP5 not closed).

The layer spends no guard:
\[
R_{\rm glob}^{(v490\text{-}r80,+)}=(0,1,1,1,1),
\qquad
\Gamma_{\rm guard}^{v490\text{-}r80}=(1,1,1,1,1),
\qquad
\beta_4=1/256\ \text{unspent}.
\]
It does not imply exact finite \(C5\) zero closure, exact all-readout terminality,
\(C6/T2\), a global photon-clock/RH theorem, or global all-readout closure.
